\section*{8. 任意多个基本形式所成系统的不变量的整有理表示}
\begin{center}
Math. Ann. 77 (1916)，第 93--102 页
\end{center}

\begin{center}
{\Large 任意多个基本形式所成系统的不变量的整有理表示。}\par
\vspace{1.2em}
作者\par
\vspace{0.8em}
{\scshape Emmy Noether}（埃尔朗根）。
\end{center}

在为 Schwarz 纪念文集撰写的《任意多个基本形式所成系统的不变量》一文末尾，D. Hilbert 提出猜想：\emph{这些不变量可以用系统 $(J,PJ)$ 的有限多个不变量作整有理表示}；这里，$J$ 表示基本形式中线性无关者的完全不变量系统，而不变量 $PJ$ 则由 $J$ 经极化过程得到。

下文的第一节将说明，这一猜想确实成立，并且它是下述约化定理的一个直接推论：\emph{对于任意多个变量组、每组各含 $n$ 个变量的任一形式，都可以把它表示为若干形式的极化式之和；这些形式仅含 $n$ 个固定的变量组，而且其本身均由原形式经极化过程导出}。这个约化定理最先由 F. Mertens 表述；它是 A. Capelli、F. Mertens 和 J. Deruyts 所给出的 Clebsch--Gordan 级数展开之推广的一个特殊情形。该推广适用于任意多个变量组、每组各含 $n$ 个变量的形式，并且是通过对微分过程作形式变换而证明的。\NoetherSrcNote{*)}{F. Mertens：《Über eine Formel der Determinantentheorie》（公式 5），Sitzb. d. Ak. d. Wiss. Wien，第 91 卷，第 II 部（1885）；更详尽的论述（包括应用）见《Über ganze Funktionen von $m$ Systemen von je $n$ Unbestimmten》，Monatsh. f. Math. u. Phys.，第 4 年卷（1893）。关于 Capelli 的若干证明（始于 1882 年），例如可参见 Math. Ann. 37，或其摹写印行的讲义《Lezioni sulla teoria delle forme algebriche》（1902）。J. Deruyts：\emph{Essai d'une theorie générale des formes algébriques}，Mém. d. l. Soc. Roy. des Sciences de Liège (2) 17（1892）。}

第二节还将给出约化定理的一个更具概念性的证明，其中把这个问题视为形式的线性族的等价问题。这一观点以及所用论证的一个实质部分，取自 E. Fischer 的论文《论代数模系统》\NoetherSrcNote{*)}{E. Fischer：Über algebraische Modulsysteme und lineare homogene partielle Differentialgleichungen mit konstanten Koeffizienten，\S\ 1--4，J. f. M. 140（1911）。}及其后续未发表的定理；E. Fischer 友好地准许我使用这些结果。

最后，第三节还将说明，相应的论证如何导出最一般的级数展开。

\subsection*{I.}
设 $\theta$ 是一个形式，它对于下列 $N>n$ 个变量组中的每一组
\[
  A_1,A_2,\ldots,A_N
\]
都是齐次的；一般地，变量组 $A_k$ 由 $n$ 个元素
\[
  A_k^{(1)},A_k^{(2)},\ldots,A_k^{(n)}
\]
组成。$\theta$ 的系数可以是不定元，也可以是某个给定有理性域中的量。再令 $P$ 表示极化过程的一个整有理函数——\emph{其系数为有理数}——其中
\[
  P_{hk}=\left(A_h\frac{\partial}{\partial A_k}\right)
  =A_h^{(1)}\frac{\partial}{\partial A_k^{(1)}}+
   A_h^{(2)}\frac{\partial}{\partial A_k^{(2)}}+
   \cdots+
   A_h^{(n)}\frac{\partial}{\partial A_k^{(n)}},
\]
而乘积 $P_{hk}P_{ij}\theta$ 应理解为把 $P_{hk}$ 作用于 $P_{ij}\theta$。

于是，上述\emph{约化定理}由恒等式
\begin{equation}
  \theta=\sum PZ,
\tag{1}
\end{equation}
给出，其中各形式 $Z$ 只含变量组
\[
  A_1,A_2,\ldots,A_n
\]
并且均由 $\theta$ 经极化过程 $P$ 导出。

现在考虑基本形式系统 $(F')$，它由 $N$ 个次数相同、变量数也相同的形式
\[
  F_1,F_2,\ldots,F_N,
\]
组成，并分别把它们的系数取作 $N$ 个变量组
\[
  A_1,A_2,\ldots,A_N
\]
。以 $(J)$ 表示由有限多个不变量组成的 $F_1,\ldots,F_n$ 的完全系统，使得 $F_1,\ldots,F_n$ 的每一个不变量都能用 $(J)$ 中的不变量作整有理表示。\emph{这样，所要证明的 Hilbert 猜想便是：对于同时不变量 $S$（属于 $F_1,\ldots,F_N$），有限多个不变量 $(J,PJ)$ 构成一个完全系统；这里，$PJ$ 指一切经过极化过程 $P$ 而可由 $J$ 导出的不变量。}

事实上，每一个有理数系数的同时不变量 $S$（属于 $(F')$），都是一个对于 $N$ 个变量组 $A_1,\ldots,A_N$ 中的每一组均齐次且具有有理数系数的形式，因而可以对它应用恒等式 (1)。众所周知，把极化过程 $P$ 施行于 $S$ 仍会产生不变量，所以各形式 $Z$ 也仍是不变量；而且它们是 $F_1,
\ldots,F_n$ 的同时不变量，因为它们只含变量组 $A_1,
\ldots,A_n$。按照假设，它们因而是 $(J)$ 中各不变量的整有理函数。恒等式 (1) 遂化为
\[
  S=\sum PY,
\]
其中 $Y$ 表示 $(J)$ 中不变量的幂乘积。可是，把过程 $P$ 实际施行于 $Y$，按照 $P$ 的定义，就是有限次地相继施行简单极化运算 $P_{hk}$；根据乘积微分法则，其结果是由 $(J)$ 和 $(PJ)$ 中的不变量构成的乘积之和。把 $P_{hk}$ 施行于由 $(J)$ 和 $(PJ)$ 中不变量构成的乘积时，情形同样如此；因此，$PY$ 也只给出 $(J,PJ)$ 的整有理组合。\emph{至此已经证明，所有不变量 $S$ 都可以用 $(J,PJ)$ 作整有理表示。}
